**Claim.** \(\sqrt{7}\) is irrational.

**Proof.** We prove the claim by contradiction.

First, we record the following elementary fact.

**Lemma.** If \(n\in \mathbb Z\) and \(7\mid n^2\), then \(7\mid n\).

**Proof of lemma.** By the division algorithm, every integer \(n\) can be written as

\[
n=7q+r
\]

for some \(q\in \mathbb Z\) and some \(r\in \{0,1,2,3,4,5,6\}\). Then

\[
n^2 \equiv r^2 \pmod 7.
\]

Checking the possible residues,

\[
0^2\equiv 0,\quad 1^2\equiv 1,\quad 2^2\equiv 4,\quad 3^2\equiv 2,
\]
\[
4^2\equiv 2,\quad 5^2\equiv 4,\quad 6^2\equiv 1 \pmod 7.
\]

Thus \(r^2\equiv 0\pmod 7\) only when \(r=0\). Therefore, if \(7\mid n^2\), then \(r=0\), so \(n=7q\), and hence \(7\mid n\). \(\square\)

Now assume, for contradiction, that \(\sqrt{7}\) is rational. By definition of rationality, there exist integers \(m,n\) with \(n\neq 0\) such that

\[
\sqrt{7}=\frac{m}{n}.
\]

We may assume \(n>0\) by replacing both \(m\) and \(n\) by \(-m\) and \(-n\) if necessary.

Let

\[
g=\gcd(m,n).
\]

Since \(g\mid m\) and \(g\mid n\), there exist integers \(a,b\) such that

\[
m=ga,\qquad n=gb.
\]

Then \(b\neq 0\), and

\[
\sqrt{7}=\frac{m}{n}=\frac{ga}{gb}=\frac{a}{b}.
\]

Moreover, \(\gcd(a,b)=1\). Indeed, if some positive integer \(d\) divides both \(a\) and \(b\), then \(dg\) divides both \(m=ga\) and \(n=gb\). Since \(g=\gcd(m,n)\) is the greatest positive common divisor of \(m\) and \(n\), we must have

\[
dg\le g.
\]

Because \(g>0\), this implies \(d\le 1\). Hence \(d=1\), so \(\gcd(a,b)=1\). Thus we have written \(\sqrt{7}\) in lowest terms:

\[
\sqrt{7}=\frac{a}{b},\qquad \gcd(a,b)=1.
\]

Squaring both sides gives

\[
7=\frac{a^2}{b^2}.
\]

Multiplying by \(b^2\), we obtain

\[
a^2=7b^2.
\]

Therefore \(7\mid a^2\). By the lemma, \(7\mid a\). Hence there exists \(k\in \mathbb Z\) such that

\[
a=7k.
\]

Substituting into \(a^2=7b^2\), we get

\[
(7k)^2=7b^2.
\]

So

\[
49k^2=7b^2.
\]

Dividing by \(7\), we obtain

\[
7k^2=b^2.
\]

Thus \(7\mid b^2\). By the lemma again, \(7\mid b\).

We have shown that \(7\mid a\) and \(7\mid b\). Therefore \(7\) is a positive common divisor of \(a\) and \(b\). Since \(7>1\), this implies

\[
\gcd(a,b)\ge 7,
\]

contradicting the earlier conclusion that \(\gcd(a,b)=1\).

Hence our assumption that \(\sqrt{7}\) is rational must be false. Therefore,

\[
\boxed{\sqrt{7}\text{ is irrational}.}
\]